\documentclass[11pt,a4paper]{article}
\usepackage[T1]{fontenc}\usepackage{lmodern}\usepackage[scaled=0.92]{helvet}
\usepackage[margin=2.2cm]{geometry}
\usepackage{amsmath,amssymb,booktabs,array,tabularx,graphicx,xcolor,enumitem,tikz}
\usetikzlibrary{positioning,arrows.meta,calc,fit,backgrounds,shapes.geometric}
\definecolor{mechblue}{HTML}{2A78D6}\definecolor{advorange}{HTML}{C9541F}\definecolor{okgreen}{HTML}{0E7C4F}
\usepackage{xr-hyper}\externaldocument{../manuscript/main}
\usepackage[hidelinks]{hyperref}
\usepackage[most]{tcolorbox}
\usepackage{titlesec}
\titleformat{\section}{\sffamily\Large\bfseries\color{mechblue!80!black}}{\thesection}{0.6em}{}
\titleformat{\subsection}{\sffamily\large\bfseries}{\thesubsection}{0.5em}{}
\setlength{\parindent}{0pt}\setlength{\parskip}{0.55em}
\setlist{itemsep=0.15em,topsep=0.3em}
\newtcolorbox{analogy}[1][]{colback=okgreen!6,colframe=okgreen!70!black,fonttitle=\sffamily\bfseries,title={In everyday terms#1},boxrule=0.6pt,arc=2pt}
\newtcolorbox{keypoint}[1][]{colback=mechblue!6,colframe=mechblue!80!black,fonttitle=\sffamily\bfseries,title={#1},boxrule=0.6pt,arc=2pt}
\newtcolorbox{caution}[1][]{colback=advorange!6,colframe=advorange!80!black,fonttitle=\sffamily\bfseries,title={#1},boxrule=0.6pt,arc=2pt}
\newcommand{\mech}[1]{\textsf{\textcolor{mechblue!80!black}{#1}}}
\newcommand{\adv}[1]{\textsf{\textcolor{advorange!85!black}{#1}}}
\newcommand{\cmd}[1]{\texttt{\small #1}}

\title{\sffamily\bfseries What Can Provenance Prove?\\[3pt]
\large\mdseries A plain-language study guide to the paper (manuscript v12)}
\author{Waqas Ahmed \and Ammar Hassan \and Attiq Ahmed \and Waseem Iqbal}
\date{21 November 2026}
\begin{document}
\sloppy
\maketitle

\begin{keypoint}{How to read this guide}
No background in machine learning security is assumed. Each idea is given in plain words,
then as an everyday comparison (green boxes), then with the paper's own term so it can be
found in the manuscript. Orange boxes mark what the paper does \emph{not} claim. Reading time:
about 30 minutes.
\end{keypoint}

\tableofcontents
\newpage

\section{The paper in one minute}
Large language models (LLMs) learn from huge collections of text. If someone slips bad
examples into that text (\emph{data poisoning}), the model can learn hidden bad behaviour,
such as a \emph{backdoor}: it behaves normally until it sees a secret trigger word.

Many researchers propose \emph{provenance} as a defence: keep signed, tamper-evident records
of where every piece of training data came from and what happened to it, sometimes on a
blockchain. Our paper asks a simple question that is rarely answered precisely:

\begin{center}\sffamily\large
When all the provenance checks pass, what has actually been proved?
\end{center}

We answer it exactly, with mathematics checked by computer, and then test the answer with
real software and real GPU training. The short answer: provenance can prove \emph{who}
supplied the data and \emph{that} the model was trained on it, under stated conditions. It can
never prove the data was \emph{harmless}.

\begin{analogy}
Think of a restaurant. Training data are the ingredients, the model is the dish, and the
auditor is a food inspector. Provenance is the paperwork: delivery notes, supplier
certificates, kitchen logs. Paperwork can prove which approved supplier delivered which
ingredient and that the dish was cooked from those ingredients. It cannot prove that an
approved supplier did not deliver spoiled food with perfectly valid paperwork.
\end{analogy}

\section{Key words}
\begin{center}\small
\begin{tabularx}{\textwidth}{@{}l>{\raggedright\arraybackslash}X@{}}
\toprule
Term & Meaning\\
\midrule
Training data & The examples a model learns from.\\
Pipeline & The steps data pass through: sources $\to$ (generator) $\to$ transform (filtering) $\to$ capture (writing the record) $\to$ executor (training).\\
Provenance & Records of where data came from and what was done to it.\\
Poisoning, backdoor & Deliberately harmful training examples; a backdoor is a hidden behaviour switched on by a trigger.\\
Auditor & Whoever checks the records before trusting a model (a regulator, customer, or internal team).\\
Mechanism & A technical safeguard, e.g.\ signed receipts or trusted hardware. We use seven.\\
Adversary & The attacker: which parts of the pipeline are dishonest. We use seven ``capabilities''.\\
Assurance goal & What the auditor wants to be sure of. We use six.\\
Binding & Proof that the released model really came from training on the recorded data.\\
Seed & A number that controls the random choices made during training.\\
Grinding & Training many times with different seeds and keeping the result you like.\\
Detector & A tool that looks for poisoned examples.\\
Recall, precision & Recall: share of poison the detector finds. Precision: share of flagged items that really are poison.\\
AUROC & One number for how well a detector separates poison from clean data: 1 = perfect, 0.5 = guessing.\\
TEE / enclave & Trusted hardware that runs code so even the machine's owner cannot tamper with it.\\
Model checking (Alloy) & A computer tool that checks a mathematical model against every possible case up to a size.\\
\bottomrule
\end{tabularx}
\end{center}

\section{The setting}
Figure~\ref{fig:sg-pipe} shows the pipeline we study. Grey boxes are the stages. Blue labels
show where each safeguard (mechanism) acts; orange labels show what each kind of attacker can
corrupt.

\begin{figure}[h]
\centering
\resizebox{0.82\textwidth}{!}{\input{../manuscript/fig/fig_pipeline}}
\caption{The pipeline, the seven mechanisms (blue) and the seven adversary capabilities
(orange). Same figure as Fig.~\ref{fig:pipeline} of the paper.}\label{fig:sg-pipe}
\end{figure}

\begin{center}\small
\begin{tabularx}{\textwidth}{@{}l>{\raggedright}X>{\raggedright\arraybackslash}X@{}}
\toprule
Mechanism & What it does & Restaurant version\\
\midrule
\mech{WitSrc} & each supplier signs a receipt for each item & supplier's signed delivery note\\
\mech{WitIn} & an independent observer records what was fed to training & an inspector watching what goes into the pot\\
\mech{AttTrans}, \mech{AttCap} & trusted hardware runs the filtering and the record-keeping & sealed, tamper-proof prep station and weighing scale\\
\mech{BindFix}, \mech{BindGrind} & proof the model came from the recorded data, with the seed fixed by the checker (Fix) or chosen by the trainer (Grind) & re-cooking the recipe to check you get the same dish\\
\mech{Detector} & looks for poisoned items & a taster\\
\bottomrule
\end{tabularx}
\end{center}

\section{How we did it}
\begin{figure}[h]
\centering
\resizebox{\textwidth}{!}{\input{../manuscript/fig/fig_overview}}
\caption{The six steps of the work (Fig.~\ref{fig:overview} of the paper).}
\end{figure}

\begin{enumerate}
\item \textbf{Systematise.} We read the literature and reduced published systems to seven basic
  mechanisms.
\item \textbf{Model.} We wrote the pipeline, attacker and goals as a precise mathematical model.
\item \textbf{Decide.} A computer checked every combination: 6 attacker types $\times$ 2 pipeline
  types $\times$ 128 mechanism combinations $\times$ 6 goals = 9,216 cases, by two independent
  methods that agree on every case; then all 128 possible attacker types.
\item \textbf{Prove.} We turned the patterns into six theorems and proved them.
\item \textbf{Validate.} We built a working prototype and attacked it 3.5 million times, then
  repeated key experiments with real language models on GPUs.
\item \textbf{Advise.} We turned the results into a checklist for auditors.
\end{enumerate}

\begin{analogy}[: why checking one item is enough]
The \emph{separability lemma} says that if a rule can be broken at all, it can be broken by a
single item. So instead of checking every possible dataset (infinitely many), it is enough to
check every possible single item, which is a small, finite job. This is why the results are
exact rather than estimates.
\end{analogy}

\section{The six results in plain words}

\subsection*{1. Binding is needed exactly when the trainer is not trusted}
If the organisation that trains the model might cheat, nothing else (signatures, witnesses,
trusted hardware earlier in the pipeline) can replace a proof that the model came from the
recorded data. If the trainer is trusted, that expensive proof is unnecessary.
\begin{analogy}
If you do not trust the cook, delivery notes do not tell you what went into the dish; you need
to watch the cooking or re-cook it yourself. If the cook is trusted, the notes are enough.
\end{analogy}

\subsection*{2. Two independent routes to tracing data back to approved sources}
Data can be traced to approved sources either through signed supplier receipts or through
trusted hardware at the filtering and recording stages. Each route fails when its own trust
anchor fails. A related finding (Corollary~1): the pipeline must declare honestly which steps
create synthetic data; otherwise smuggled data can be labelled ``synthetic'' and escape the
supplier checks.
\begin{analogy}
Smuggled ingredients labelled ``made in-house'' are never checked against supplier notes. So
the kitchen must declare in advance what it actually makes in-house.
\end{analogy}

\subsection*{3 and 4. Choosing among training runs leaks control (``grinding'')}
Training involves random choices. If the trainer may pick the random seed, it can train many
times and release the run that behaves as it wants, and the checks still pass. With $T$ runs it
can set about $\log_2 T$ yes/no behaviours of its choice: 128 runs give 7, 512 runs give 9.
The fix is to let the checker fix the seed \emph{after} the data are committed.
\begin{analogy}
A cook who may roll the dice himself can roll 128 times and keep the result he likes; your
re-cooking with his dice result will match perfectly. If you roll the dice after the
ingredients are sealed, he cannot choose.
\end{analogy}

\subsection*{5. Provenance can never prove the model is poison-free}
An approved source can supply poison with perfectly valid paperwork. The strongest honest claim
is conditional: ``every poisoned item that was used was examined by the detector and missed''.
That is only as good as the detector.
\begin{analogy}
Valid paperwork for spoiled food is still valid paperwork. Only a taster can catch spoiled
food, and tasters miss some.
\end{analogy}

\subsection*{6. Use two independent ways to make the record complete}
Trusted hardware at the recording step and an independent observer of the training input
each make sure nothing is left out of the record. They fail for different reasons, so using
both protects against the failure of either.

\begin{keypoint}{Where a blockchain fits}
A blockchain or public log stops anyone from secretly rewriting or suppressing entries. It
does \emph{not} make the entries true, complete or poison-free. It is useful only when
suppression (censorship) is a real risk.
\end{keypoint}

\section{What the experiments showed}
\begin{description}[leftmargin=0pt]
\item[The model matches reality.] In 3.5 million attacked runs of the prototype, no accepted run
  broke a guarantee the model gives, and every failure the model predicted was demonstrated by
  a real attack.
\item[Grinding works exactly as predicted,] on a small model (9 chosen behaviours with 512 runs)
  and on real language models (7 with 128 DistilBERT runs; 5 with 32 Qwen runs).
\item[Re-training reproduces the model] bit for bit on the same machine (G2) but not on a
  different GPU with newer software (G5). Checking by re-training needs the same hardware and
  software.
\item[Poison passes provenance.] Poisoned examples from an approved source carried valid
  signatures. At 2\% poison the backdoor worked, while the detector found only about a quarter
  of the poison (G3).
\item[Detectors fail when it matters most (G5).] Figure~\ref{fig:sg-g5} shows each training run:
  detectors separate poison well when the attack is not working, and are no better than
  guessing once the backdoor works every time.
\end{description}

\begin{figure}[h]
\centering
\includegraphics[width=0.62\textwidth]{../manuscript/fig/g5_detection.pdf}
\caption{G5: how well two detectors separate poisoned from clean examples (up = better;
dashed line = guessing) against how well the backdoor works (right = attack succeeds). One
point per training run.}\label{fig:sg-g5}
\end{figure}

\begin{analogy}
The taster notices a little spoiled food that has not yet changed the dish. Once there is
enough to really change the dish, the spoiled food tastes ``normal'' to the taster.
\end{analogy}

\section{The auditor's checklist}
\begin{figure}[h]
\centering
\resizebox{0.75\textwidth}{!}{\input{../manuscript/fig/fig_auditor}}
\caption{Four questions an auditor should ask (Fig.~\ref{fig:auditor} of the paper).}
\end{figure}

\section{What the paper does not claim}
\begin{caution}{Limits}
\begin{itemize}
\item It does not claim provenance or blockchain prevents poisoning.
\item The trusted hardware in the prototype is simulated, not real.
\item Computer checking is exact for the model at the stated size; the model is a simplified
  picture of real pipelines.
\item Poisoning was tested on one model family, one dataset and one trigger.
\item The experiments cover fine-tuning, not training a large model from scratch.
\item The literature screening was done by an AI assistant and still needs human screening.
\end{itemize}
\end{caution}

\section{Questions you may be asked}
\begin{description}[leftmargin=0pt]
\item[What is new compared with existing surveys?] Surveys list techniques. We decide exactly which
  combinations guarantee which goals against which attackers, prove it, and test it.
\item[Isn't the grinding attack known?] Tan et al.\ (CCS 2025) showed it for zero-knowledge proofs
  of training. We give a general limit ($\log_2 T$ bits) that also covers non-convex training
  such as LLM fine-tuning, and measure it on real models.
\item[Why trust the computer checks?] Two independent methods agree on all 9,216 cases, all
  128 attacker types were also checked, and the prototype matched the model in 3.5 million runs.
\item[Why did G5 use a different GPU?] The original laptop was unavailable. The paper says so,
  and this produced a useful finding: re-training on a different machine does not reproduce the
  model exactly.
\item[Does a blockchain help?] Only against suppression of records, not against false or
  poisoned data (Proposition~2).
\end{description}

\section{What remains before submission}
\begin{itemize}
\item Database searches (Scopus, Web of Science, IEEE Xplore, ACM), full-text screening and a
  second human screener.
\item The supervisor's check of the proofs of Theorems~1 and~4.
\item A short statement of what the authors verified themselves, and a check of the AI-use
  statement against IEEE policy.
\item Departments, postal addresses and e-mail addresses.
\item Optional: a poisoning experiment on a second model family (Qwen2.5-0.5B with LoRA).
\end{itemize}
\end{document}
